\documentclass{article}

\usepackage{amsmath,amssymb}
\usepackage{xurl}
\usepackage[hidelinks]{hyperref}
\setlength{\emergencystretch}{2em}

\input{command}

\title{Schur Complement TFAE --- Missing Left-Support Condition}
\author{}
\date{2026-08-21}

\begin{document}
\maketitle
\gapnote{false-source}{resolved}

\section{Source statement}

Let $P$ and $R$ be positive semidefinite complex matrices, and let $Q$ be a
complex matrix of compatible dimensions. Theorem~5.2 of
Ref.~\cite{Wolf2012Quantum} states that the following conditions are
equivalent:
\begin{enumerate}
    \item The block matrix
    \begin{align}
        B&=
        \begin{pmatrix}
            P & Q\\
            Q^\dagger & R
        \end{pmatrix}
    \label{eq:wolf_schur_complement_tfae_block_matrix}
    \end{align}
    is positive semidefinite.

    \item The right-support condition
    \begin{align}
        \ker(R)&\subseteq\ker(Q)
    \label{eq:wolf_schur_complement_tfae_right_support}
    \end{align}
    holds, and
    \begin{align}
        P&\geq QR^+Q^\dagger,
    \label{eq:wolf_schur_complement_tfae_schur_inequality}
    \end{align}
    where $R^+$ is the pseudoinverse, equal to the inverse on the support of
    $R$.

    \item The condition in
    (\ref{eq:wolf_schur_complement_tfae_right_support}) holds, and
    \begin{align}
        \left\lVert P^{-1/2}QR^{-1/2}\right\rVert_\infty
            &\leq1,
    \label{eq:wolf_schur_complement_tfae_printed_contraction}
    \end{align}
    where the inverse square roots vanish on the corresponding kernels.
\end{enumerate}

The third condition in Theorem~5.2 of Ref.~\cite{Wolf2012Quantum} is false
when $P$ is singular.

\section{Counterexample and correction}

Take one-dimensional matrices
\begin{align}
    P&=0,
    \label{eq:wolf_schur_complement_tfae_counterexample_p}\\
    Q&=1,
    \label{eq:wolf_schur_complement_tfae_counterexample_q}\\
    R&=1.
    \label{eq:wolf_schur_complement_tfae_counterexample_r}
\end{align}
Then $\ker(R)\subseteq\ker(Q)$, and the support inverse gives
\begin{align}
    P^{-1/2}QR^{-1/2}
        &=0,
    \label{eq:wolf_schur_complement_tfae_counterexample_contraction}
\end{align}
so the printed contraction condition
(\ref{eq:wolf_schur_complement_tfae_printed_contraction}) holds. However, for
$x=(1,-1)^{\mathsf T}$,
\begin{align}
    \begin{pmatrix}1&-1\end{pmatrix}
    \begin{pmatrix}0&1\\1&1\end{pmatrix}
    \begin{pmatrix}1\\-1\end{pmatrix}
        &=-1.
    \label{eq:wolf_schur_complement_tfae_counterexample_quadratic_form}
\end{align}
Thus the block matrix is not positive semidefinite.

The missing hypothesis is the left-support condition
\begin{align}
    \ker(P)&\subseteq\ker(Q^\dagger).
    \label{eq:wolf_schur_complement_tfae_left_support}
\end{align}

\section{The Schur-complement equivalence}

The formal development defines the block matrix
(\ref{eq:wolf_schur_complement_tfae_block_matrix}), proves its
self-adjointness for Hermitian diagonal blocks, and defines the pseudoinverse
Schur complement
\begin{align}
    S&=P-QR^+Q^\dagger.
    \label{eq:wolf_schur_complement_tfae_schur_complement}
\end{align}
The corresponding declarations are
\leanid{SchurComplement.blockMatrix},
\leanid{SchurComplement.blockMatrix_isHermitian}, and
\leanid{SchurComplement.schurComplement} in
\doclink{QICLean/Channel/Schwarz/SchurComplement}.

Conditions one and two of Theorem~5.2 in Ref.~\cite{Wolf2012Quantum} are
equivalent under the source hypotheses:
\begin{align}
    B\geq0
        &\Longleftrightarrow
        \left(\ker(R)\subseteq\ker(Q)
            \ \text{and}\
            P-QR^+Q^\dagger\geq0\right).
    \label{eq:wolf_schur_complement_tfae_block_schur_equivalence}
\end{align}
This is
\leanid{SchurComplement.blockMatrix_posSemidef_iff}. The forward implication
combines
\leanid{SchurComplement.ker_subset_of_block_psd} with quadratic-form
minimization at
\begin{align}
    y&=-R^+Q^\dagger x.
    \label{eq:wolf_schur_complement_tfae_minimizing_vector}
\end{align}
The reverse implication uses the completed-square identity
\begin{align}
    \begin{pmatrix}x\\y\end{pmatrix}^{\!\dagger}
    \begin{pmatrix}P&Q\\Q^\dagger&R\end{pmatrix}
    \begin{pmatrix}x\\y\end{pmatrix}
        &=
        (R^+Q^\dagger x+y)^\dagger
        R(R^+Q^\dagger x+y)
    \label{eq:wolf_schur_complement_tfae_completed_square_positive_term}\\
        &\quad+
        x^\dagger(P-QR^+Q^\dagger)x.
    \label{eq:wolf_schur_complement_tfae_completed_square_schur_term}
\end{align}
This identity is formalized as
\leanid{SchurComplement.schur_complement_quadratic_form}. Its proof uses
\begin{align}
    Q\operatorname{supp}(R)&=Q,
    \label{eq:wolf_schur_complement_tfae_support_absorption}\\
    RR^+&=\operatorname{supp}(R),
    \label{eq:wolf_schur_complement_tfae_pseudoinverse_left}\\
    R^+R&=\operatorname{supp}(R),
    \label{eq:wolf_schur_complement_tfae_pseudoinverse_right}\\
    (R^+)^\dagger&=R^+.
    \label{eq:wolf_schur_complement_tfae_pseudoinverse_adjoint}
\end{align}

\section{The corrected contraction condition}

The scalar counterexample is formalized as
\leanid{SchurComplement.wolf_condition_three_not_sufficient}. To state the
corrected condition, define
\begin{align}
    K&=P^{-1/2}QR^{-1/2}.
    \label{eq:wolf_schur_complement_tfae_contraction_operator}
\end{align}
Assume both support conditions
(\ref{eq:wolf_schur_complement_tfae_right_support}) and
(\ref{eq:wolf_schur_complement_tfae_left_support}). Then
\begin{align}
    KK^\dagger
        &=
        P^{-1/2}(QR^+Q^\dagger)P^{-1/2}.
    \label{eq:wolf_schur_complement_tfae_contraction_product}
\end{align}
Consequently,
\begin{align}
    P-QR^+Q^\dagger\geq0
        &\Longleftrightarrow
        \lVert K\rVert_\infty\leq1.
    \label{eq:wolf_schur_complement_tfae_corrected_contraction_equivalence}
\end{align}
This equivalence is formalized as
\leanid{SchurComplement.schurComplement_posSemidef_iff_contraction}; its
block-matrix form is
\leanid{SchurComplement.blockMatrix_posSemidef_iff_contraction}.

\section{Verdict}

Conditions one and two of Theorem~5.2 in
Ref.~\cite{Wolf2012Quantum} are correct as printed and are formalized without
additional hypotheses. Condition three is false as printed, as
(\ref{eq:wolf_schur_complement_tfae_counterexample_p})--%
(\ref{eq:wolf_schur_complement_tfae_counterexample_quadratic_form}) show.
Adding the necessary left-support condition
(\ref{eq:wolf_schur_complement_tfae_left_support}) restores the three-way
equivalence. Thus
\href{https://github.com/LionSR/TNLean/issues/5501}{TNLean issue \#5501} is
resolved by correcting the source statement rather than formalizing the false
original assertion.

\bibliographystyle{alpha}
\bibliography{references}

\end{document}
